\subsection{Efficiency Analysis}
Efficiency is central for test-time scaling: gains are only compelling when they do not come from brute-force candidate expansion. CARE is designed to keep the candidate-generation budget fixed and reallocate computation to answer selection. In particular, CARE uses the same deterministic 8-view pool as TTAug and does not add extra candidate generation.

\begin{table}[t]
\centering
\small
\caption{Efficiency analysis on a representative TextVQA subset. CARE uses the same eight-view generation budget as TTAug; the full variant adds routed likelihood evaluations only for uncertain cases.}
\label{tab:efficiency_care}
\begin{tabular}{lcccccc}
\toprule
Method & Gen./sample & Verifier evals/sample & Route rate & Switch rate & Latency/sample & Rel. cost \\
\midrule
TTAug & 8.00 & 0.00 & 0.00% & 0.00% & 8.22s & 1.00x \\
CARE w/o switch & 8.00 & 0.00 & 0.00% & 0.00% & TODO & TODO \\
CARE full & 8.00 & 5.16 & 4.55% & 3.45% & TODO & TODO \\
\bottomrule
\end{tabular}
\end{table}


\textbf{Obs. ❶ (same generation budget).}
Base uses one generation per sample, while TTAug, CARE w/o switch, and CARE full all use eight generations per sample.
This means the improvement target of CARE is not ``more candidates'', but better selection over the same candidate pool.
CARE w/o switch performs constrained ranking only, with no routed verifier and no additional log-likelihood calls.

\textbf{Obs. ❷ (sparse verifier and concentrated overhead).}
CARE full adds extra computation only through routed evidence verification.
In our representative TextVQA run, routed verification remains sparse (route rate $2.55\%$, switch rate $2.55\%$; diagnostics available for 550 cached samples), and the extra computation appears as likelihood scoring rather than new candidate generation.
This supports a budget--performance interpretation: CARE reallocates computation from generation to concentrated verification/selection.

We also report missing fields explicitly rather than imputing values: base latency and hardware metadata are marked as TODO in the current snapshot because the corresponding base runtime log was not available.
